\chapter{Propuesta económica formal}
Esta oferta establece el precio, el calendario de facturación y las condiciones bajo las cuales LafroX SpA implementará y operará la plataforma logística de Distribuidora Puelche S.A. Las condiciones comerciales se presentan según el E-21 y se sustentan en el alcance técnico, el análisis económico-financiero y el modelo de la oferta (Bases Administrativas, formulario E-21, pp.~70--72; LafroX SpA, 2026a, 2026d).
\section{Propuesta de valor económico}
El precio base reúne los doce hitos de implementación y los 36 períodos mensuales de operación. La conciliación del precio se apoya en las hojas RESUMEN y FLUJO CAJA NETO del modelo (LafroX SpA, 2026d).
\subsection{Precio por fase}
LafroX SpA ofrece a Distribuidora Puelche S.A. la implementación en dos etapas y la operación gestionada de la solución logística durante 36 meses. El precio base cubre los servicios y productos detallados en la sección 1.3. La Etapa 1 abarca los meses 1 a 16 y la Etapa 2, los meses 13 a 21, incluido el paso a producción de cada etapa (Bases Administrativas, arts.~15 y 17, pp.~11--13).

\begin{tablalafrox}{Precio base e IVA}{tab:e1-0-4}{Y R{2.1cm} R{2.8cm} R{2.4cm}}
{\cab{Concepto} & \cab{UF} & \cab{CLP} & \cab{USD}}
Implementación · neto & 115.947,39 & 4.637.895.572 & 5.153.217 \\
Implementación · IVA & 22.030,00 & 881.200.159 & 979.111 \\
Implementación · con IVA & 137.977,39 & 5.519.095.731 & 6.132.329 \\
Operación · neto & 130.068,25 & 5.202.730.094 & 5.780.811 \\
Operación · IVA & 24.712,97 & 988.518.718 & 1.098.354 \\
Operación · con IVA & 154.781,22 & 6.191.248.812 & 6.879.165 \\
Proyecto · neto & 246.015,64 & 9.840.625.666 & 10.934.029 \\
Proyecto · IVA & 46.742,97 & 1.869.718.877 & 2.077.465 \\
Proyecto · con IVA & 292.758,61 & 11.710.344.543 & 13.011.494 \\
\end{tablalafrox}

La Tabla~\ref{tab:e1-0-4} permite conciliar el precio total con la suma de la implementación y las 36 cuotas base.

Al 25 de noviembre de 2026 se aplican las equivalencias del E-24: 1 UF = 40.000 CLP y 1 USD = 900 CLP. Se calcula el 19\,\% de IVA antes de redondear a dos decimales en UF y a la unidad en CLP y USD (Bases Administrativas, art.~9, p.~7, y E-24, p.~73).

\subsection{Precio por etapa de implementación}
El E-25 reparte el neto de implementación en 60\,\% para H1--H7 y 40\,\% para H8--H12 (Bases Administrativas, pp.~74--75).
\begin{tablalafrox}{Precio por etapa}{tab:e1-1-23}{Y R{2.1cm} R{2.8cm} R{2.4cm}}
{\cab{Concepto} & \cab{UF} & \cab{CLP} & \cab{USD}}
Etapa 1 · neto & 69.568,43 & 2.782.737.343 & 3.091.930 \\
Etapa 1 · IVA & 13.218,00 & 528.720.095 & 587.467 \\
Etapa 1 · con IVA & 82.786,44 & 3.311.457.438 & 3.679.397 \\
Etapa 2 · neto & 46.378,96 & 1.855.158.229 & 2.061.287 \\
Etapa 2 · IVA & 8.812,00 & 352.480.063 & 391.645 \\
Etapa 2 · con IVA & 55.190,96 & 2.207.638.292 & 2.452.931 \\
\end{tablalafrox}

Los subtotales de la Tabla~\ref{tab:e1-1-23} permiten conciliar los pagos de los hitos con el precio de cada etapa.

\subsection{Vigencia y reajustabilidad}
La oferta rige por 90 días corridos desde su presentación. La garantía de seriedad rige por un mínimo de 150 días corridos desde la entrega, con las renovaciones del artículo 35.3 (Bases Administrativas, p.~23).

El precio de implementación es firme. La tarifa de operación es fija en UF durante 36 meses, sin reajuste por IPC.

\section{Estructura de pagos detallada}
La implementación se factura contra hitos aceptados y la operación, a mes vencido. Cada calendario tiene su propio hecho que habilita el cobro y sus importes se concilian por separado en el modelo (Bases Administrativas, formulario E-25, pp.~74--75; LafroX SpA, 2026d, hoja FLUJO CAJA NETO).
\subsection{Hitos de implementación}
Los porcentajes del E-25 se aplican al precio neto total de implementación. Cada hito se factura con IVA tras la aceptación de los entregables, acreditada mediante acta suscrita por la Contraparte Técnica del cliente conforme al art. 18.1.

\begin{tablalafrox}{Hitos E-25 y montos netos}{tab:e1-2-35}{L{1.45cm} C{0.65cm} C{1.0cm} R{1.8cm} R{2.6cm} R{2.35cm}}
{\cab{Hito} & \cab{Mes} & \cab{\%} & \cab{UF} & \cab{CLP} & \cab{USD}}
H1 & 2 & 8 \% & 9.275,79 & 371.031.646 & 412.257 \\
H2 & 4 & 7 \% & 8.116,32 & 324.652.690 & 360.725 \\
H3 & 6 & 10 \% & 11.594,74 & 463.789.557 & 515.322 \\
H4 & 10 & 10 \% & 11.594,74 & 463.789.557 & 515.322 \\
H5 & 12 & 10 \% & 11.594,74 & 463.789.557 & 515.322 \\
H6 & 13 & 5 \% & 5.797,37 & 231.894.779 & 257.661 \\
H7 & 16 & 10 \% & 11.594,74 & 463.789.557 & 515.322 \\
H8 & 14 & 5 \% & 5.797,37 & 231.894.779 & 257.661 \\
H9 & 17 & 10 \% & 11.594,74 & 463.789.557 & 515.322 \\
H10 & 18 & 10 \% & 11.594,74 & 463.789.557 & 515.322 \\
H11 & 19 & 5 \% & 5.797,37 & 231.894.779 & 257.661 \\
H12 & 21 & 10 \% & 11.594,74 & 463.789.557 & 515.322 \\
Etapa 1 &  & 60 \% & 69.568,43 & 2.782.737.343 & 3.091.930 \\
Etapa 2 &  & 40 \% & 46.378,96 & 1.855.158.229 & 2.061.287 \\
Total &  & 100 \% & 115.947,39 & 4.637.895.572 & 5.153.217 \\
\end{tablalafrox}

La Tabla~\ref{tab:e1-2-35} permite comprobar qué monto neto corresponde facturar en cada hito aceptado. El mes 21 incluye dos cobros separados, H12 y la primera cuota de operación. Los subtotales se redondean por separado.

\subsection{IVA de los hitos de implementación}
Los importes se calculan sobre los valores netos sin redondear. Los subtotales de etapa y el total de implementación se redondean de forma independiente.
\begin{tablalafrox}{IVA de los hitos de implementación en UF, CLP y USD}{tab:e1-hitos-iva}{Y R{2.1cm} R{2.8cm} R{2.4cm}}
{\cab{Hito} & \cab{UF} & \cab{CLP} & \cab{USD}}
H1 & 1.762,40 & 70.496.013 & 78.329 \\
H2 & 1.542,10 & 61.684.011 & 68.538 \\
H3 & 2.203,00 & 88.120.016 & 97.911 \\
H4 & 2.203,00 & 88.120.016 & 97.911 \\
H5 & 2.203,00 & 88.120.016 & 97.911 \\
H6 & 1.101,50 & 44.060.008 & 48.956 \\
H7 & 2.203,00 & 88.120.016 & 97.911 \\
H8 & 1.101,50 & 44.060.008 & 48.956 \\
H9 & 2.203,00 & 88.120.016 & 97.911 \\
H10 & 2.203,00 & 88.120.016 & 97.911 \\
H11 & 1.101,50 & 44.060.008 & 48.956 \\
H12 & 2.203,00 & 88.120.016 & 97.911 \\
Etapa 1 & 13.218,00 & 528.720.095 & 587.467 \\
Etapa 2 & 8.812,00 & 352.480.063 & 391.645 \\
Implementación & 22.030,00 & 881.200.159 & 979.111 \\
\end{tablalafrox}

La Tabla~\ref{tab:e1-hitos-iva} identifica el impuesto que se agrega a cada hito aceptado. Permite comprobar el IVA de una factura a partir del importe neto del hito y conciliarlo con el IVA total de la implementación.

\subsection{Monto con IVA de los hitos de implementación}
Los importes se calculan sobre los valores netos sin redondear. Los subtotales de etapa y el total de implementación se redondean de forma independiente.
\begin{tablalafrox}{Monto con IVA de los hitos de implementación en UF, CLP y USD}{tab:e1-hitos-total}{Y R{2.1cm} R{2.8cm} R{2.4cm}}
{\cab{Hito} & \cab{UF} & \cab{CLP} & \cab{USD}}
H1 & 11.038,19 & 441.527.658 & 490.586 \\
H2 & 9.658,42 & 386.336.701 & 429.263 \\
H3 & 13.797,74 & 551.909.573 & 613.233 \\
H4 & 13.797,74 & 551.909.573 & 613.233 \\
H5 & 13.797,74 & 551.909.573 & 613.233 \\
H6 & 6.898,87 & 275.954.787 & 306.616 \\
H7 & 13.797,74 & 551.909.573 & 613.233 \\
H8 & 6.898,87 & 275.954.787 & 306.616 \\
H9 & 13.797,74 & 551.909.573 & 613.233 \\
H10 & 13.797,74 & 551.909.573 & 613.233 \\
H11 & 6.898,87 & 275.954.787 & 306.616 \\
H12 & 13.797,74 & 551.909.573 & 613.233 \\
Etapa 1 & 82.786,44 & 3.311.457.438 & 3.679.397 \\
Etapa 2 & 55.190,96 & 2.207.638.292 & 2.452.931 \\
Implementación & 137.977,39 & 5.519.095.731 & 6.132.329 \\
\end{tablalafrox}

La Tabla~\ref{tab:e1-hitos-total} muestra cuánto se factura por cada hito una vez sumado el IVA. Sus subtotales permiten estimar el desembolso por etapa y cotejarlo con el precio de implementación con impuesto.


\subsection{Entregables y criterios de aceptación}
La entrega se acompaña del artefacto, evidencia objetiva de verificación, trazabilidad a los requisitos y registro de observaciones resueltas. Puelche dispone de diez días hábiles para pronunciarse y LafroX de diez días hábiles para subsanar las observaciones. La segunda presentación con observaciones de la misma naturaleza se considera atraso imputable a LafroX. La firma del acta habilita la facturación (Bases Administrativas, art.~18, p.~13; LafroX SpA, 2026e).
\begin{tablalafrox}{Conformidad requerida por hito}{tab:e1-3-56}{L{1.4cm} Y}
{\cab{Hito} & \cab{Entregable y conformidad requerida}}
H1 & Línea base de alcance y matriz de trazabilidad de requerimientos aprobadas. \\
H2 & Documento de arquitectura, plan de seguridad y modelo de datos aprobados. \\
H3 & Infraestructura híbrida entregada y ambientes de desarrollo, QA, preproducción y producción habilitados con observabilidad. \\
H4 & Software de la Etapa 1 entregado para pruebas con QA superado y evidencia de cobertura. \\
H5 & Etapa 1 certificada con pruebas de aceptación de usuarios, carga, resiliencia y seguridad ofensiva aprobadas. \\
H6 & Marcha blanca de la Etapa 1 iniciada con plan de reversión activo y usuarios capacitados. \\
H7 & Etapa 1 en producción con cierre conforme de la marcha blanca. \\
H8 & Línea base de alcance y diseño detallado de la Etapa 2 aprobados. \\
H9 & Software de la Etapa 2 entregado para pruebas con QA superado. \\
H10 & Etapa 2 certificada y desarrollo cerrado. \\
H11 & Marcha blanca de la Etapa 2 iniciada en convivencia con la Etapa 1 en producción. \\
H12 & Etapa 2 en producción con aceptación final de la implementación. \\
\end{tablalafrox}

La Tabla~\ref{tab:e1-3-56} vincula el código de cada hito de pago con la evidencia que debe revisar Puelche. Para H7 y H12, la aceptación exige además comprobar las condiciones de cierre de marcha blanca que se detallan a continuación.

\subsection{Cierre de marcha blanca}
Para H7 y H12, la puesta en producción exige que se cumplan en conjunto las condiciones de cierre del artículo 17.3. Durante al menos las cuatro últimas semanas debe alcanzarse el volumen real comprometido y cumplirse de forma sostenida la disponibilidad y los tiempos de respuesta acordados. Al cierre se verifica que no existan incidentes abiertos críticos o altos atribuibles a la solución, que la conciliación con el registro vigente carezca de diferencias sin explicar y que el personal de Puelche esté capacitado y certificado conforme al plan. La Contraparte Técnica formaliza la aceptación mediante acta. Si faltan condiciones al término del período, LafroX extiende la marcha blanca a su costo, conserva las fechas contractuales siguientes y queda sujeta a las multas por atraso (Bases Administrativas, art.~17.3, p.~13; LafroX SpA, 2026e).

\subsection{Facturación mensual de la operación}
La tarifa base cubre alojamiento e infraestructura operada, conectividad, licencias y servicios de terceros incluidos, mesa de servicio, monitoreo, seguridad y mantención de la solución. Cada uno de los 36 períodos se factura a mes vencido, con pago dentro de los primeros días del mes siguiente y con los descuentos por incumplimiento de niveles de servicio que correspondan. Desde el mes 21 se operan ambos alcances en producción. H12 corresponde a la aceptación final de la implementación y se factura por separado del servicio prestado durante ese mismo mes. La cuota base es el precio neto de operación dividido por 36, antes de IVA, descuentos o pagos variables (Bases Administrativas, art.~17.2, p.~13, y formulario E-25, pp.~74--75; LafroX SpA, 2026a, secc.~2.6; LafroX SpA, 2026d, hojas RESUMEN y FLUJO CAJA NETO).

\subsection{Importes mensuales y topes del componente variable}
La innovación INN-04 vincula una parte del pago de operación al desempeño de la solución. La cuota fija remunera los servicios incluidos durante los 36 meses de operación. A esa cuota puede sumarse un pago variable cuando se verifica una reducción del costo comparable por entrega frente a la línea base y se cumple la meta vigente de entregas completas y a tiempo (OTIF). Entre los meses 21 y 23 se construye la línea base y se ensaya el cálculo sin cobro. Del mes 24 al 56, el importe se determina mes a mes y puede alcanzar hasta 10 UF netas si la reducción comparable es de al menos 5\,\% y se cumple la meta OTIF. Si el tope se alcanza en los 33 meses elegibles, el pago variable acumulado llega a 330 UF netas. La Tabla~\ref{tab:e1-4-74} presenta por separado la cuota fija, el máximo variable de un mes y el máximo acumulado, con su desglose de IVA en UF, CLP y USD. Los dos topes corresponden al mismo componente: el acumulado representa la suma máxima de los pagos mensuales. La subsección 1.2.8 explica cómo se determina el importe efectivamente facturable (LafroX SpA, 2026a, secc.~2.7; LafroX SpA, 2026b, indicador I-04A; LafroX SpA, 2026d, hoja INNOVACIONES).

\begin{tablalafrox}{Cuota base y topes del componente variable}{tab:e1-4-74}{Y R{2.1cm} R{2.8cm} R{2.4cm}}
{\cab{Concepto} & \cab{UF} & \cab{CLP} & \cab{USD}}
Cuota base mensual neta & 3.613,01 & 144.520.280 & 160.578 \\
IVA de la cuota base & 686,47 & 27.458.853 & 30.510 \\
Cuota base mensual con IVA & 4.299,48 & 171.979.134 & 191.088 \\
Tope mensual · neto & 10,00 & 400.000 & 444 \\
Tope mensual · IVA & 1,90 & 76.000 & 84 \\
Tope mensual · con IVA & 11,90 & 476.000 & 529 \\
Tope acumulado · neto & 330,00 & 13.200.000 & 14.667 \\
Tope acumulado · IVA & 62,70 & 2.508.000 & 2.787 \\
Tope acumulado · con IVA & 392,70 & 15.708.000 & 17.453 \\
\end{tablalafrox}

La Tabla~\ref{tab:e1-4-74} permite distinguir la cuota fija de los importes máximos de INN-04. El tope acumulado corresponde a 33 meses elegibles, del mes 24 al 56. El importe que efectivamente se factura depende de la reducción de costo y del cumplimiento de la meta OTIF.

\subsection{Liquidación variable por desempeño}
En los meses 21 a 23 se construye la línea base con una canasta de entregas de pesos fijos por zona, canal, distancia y condición térmica. En esos meses se ensaya la liquidación sin generar cobro. Entre los meses 24 y 56 se utiliza la reducción comparable del costo por entrega (I-04A del T-19) y se verifica que la proporción de entregas completas y a tiempo (OTIF) cumpla la meta vigente. El T-19 define la medición de referencia y el modelo registra el tope contingente (LafroX SpA, 2026a, secc.~2.7; LafroX SpA, 2026b, indicador I-04A; LafroX SpA, 2026d, hoja INNOVACIONES).

Si se cumple la meta OTIF, el pago variable neto es de 10 UF multiplicadas por la reducción comparable dividida por 5 \%, con el factor acotado entre cero y uno. Si el OTIF queda bajo la meta, el pago variable es cero. La comparación conserva las entregas fallidas en el denominador y aplica los ajustes documentados por combustible, mezcla de productos y peak de septiembre.

Por ejemplo, con una reducción comparable de 2,5\,\% y la meta OTIF cumplida, se devengan 5,00 UF netas (200.000 CLP o 222 USD). Con una reducción del 5\,\% o superior y la meta OTIF cumplida, se alcanza el tope mensual. Una reducción nula o negativa genera un pago variable de cero.

El importe variable devengado se factura con IVA separado del precio base. Las multas se aplican según el contrato.

\section{Resumen ejecutivo de valor}
El alcance contratado comprende la plataforma híbrida, las capacidades logísticas de ambas etapas y el servicio gestionado. La oferta técnica identifica qué capacidades entran en cada etapa y qué insumos debe facilitar Puelche (LafroX SpA, 2026c, secc.~3.2; LafroX SpA, 2026e).
\subsection{Servicios y productos incluidos}
La Tabla~\ref{tab:e1-5-89} resume los servicios y productos incluidos en el precio. La Etapa 2 usa los datos operacionales consolidados en la Etapa 1. Su paso a producción incorpora el nuevo alcance al servicio gestionado desde el mes 21 (Bases Administrativas, arts.~15 y 17, pp.~11--13; LafroX SpA, 2026c, tabla~3.2).
\begin{tablalafrox}{Alcance comercial incluido}{tab:e1-5-89}{L{3.5cm} Y}
{\cab{Fase} & \cab{Servicio o producto}}
Implementación · Etapa 1 & Plataforma híbrida, ambientes, migración e integración con sistemas existentes. Comprende preventa, recepción, bodega, preparación, cross-docking, rutas, cadena de frío, entrega y acuse, rendición y retiro sanitario (F-01 a F-11). Incluye pruebas, capacitación, marcha blanca y puesta en producción. \\
Implementación · Etapa 2 & Portal, intercambio electrónico y costo de servir (F-12 a F-14), sobre la plataforma y los registros de la primera etapa. Incluye integraciones de esta etapa, pruebas, capacitación, marcha blanca en convivencia con la operación de la Etapa 1 y aceptación final. \\
Operación · Meses 21 a 56 & Operación de la plataforma, soporte N1 y especializado, NOC y SOC, gestión de incidencias, respaldo y recuperación, mantenimiento preventivo, correctivo y evolutivo, y actualizaciones según el alcance acordado. Incluye la ejecución del plan de reversibilidad al término del contrato (Bases Administrativas, art.~77.3, p.~39). \\
Innovaciones & Seguimiento de vencimiento posentrega, reproducción de incidentes de terreno, estimación de vida útil remanente, tramo variable ligado al costo de servir y hoja de negocio del almacenero (LafroX SpA, 2026b). \\
\end{tablalafrox}

La Tabla~\ref{tab:e1-5-89} delimita qué recibe Puelche en cada etapa de implementación y qué servicios quedan cubiertos durante la operación. Esta distribución permite relacionar los productos entregados con los hitos y el servicio mensual con la cuota de operación.

\subsection{Exclusiones y dependencias}
Se excluyen la sustitución del ERP, remuneraciones, punto de venta del cliente, comercio al consumidor final, administración contractual y pagos de transportistas, rediseño de la red logística, robotización y mantenimiento mecánico. La adquisición de hardware de terreno corresponde a Puelche. LafroX especifica sus cantidades y características e integra los equipos al alcance acordado. LafroX absorbe la instalación especializada incluida en Talca según el análisis financiero (LafroX SpA, 2026a, secc.~2.1; LafroX SpA, 2026c, secc.~3.2.2).

Puelche facilita acceso a datos, interfaces e instalaciones, designa contraparte y usuarios para validar y probar la solución, y disponibiliza el equipamiento de su cargo para la marcha blanca. LafroX coordina la habilitación de ambientes, las integraciones, la migración, las pruebas y la capacitación según el plan de trabajo. La Contraparte Técnica de Puelche valida los entregables y formaliza la conformidad (Bases Administrativas, art.~18, p.~13; LafroX SpA, 2026e).

\subsection{Beneficios económicos para Puelche}
La trazabilidad de lotes y vencimientos, junto con las alertas operativas, busca reducir pérdidas de mercadería y facilitar la actuación sobre productos en riesgo. La captura en terreno y la continuidad operativa buscan reducir reprocesos y tiempo detenido. La medición del costo por entrega permite revisar la atención por zona y canal con protección de la meta OTIF. Estos beneficios dependen de la adopción, la cobertura de captura y el desempeño medido durante la operación (LafroX SpA, 2026a, secc.~2.7; LafroX SpA, 2026b).

El Entregable 2 estima para INN-04 un escenario de hasta 2.046 UF de beneficio bruto acumulado en 33 meses, equivalente a 81.840.000 CLP o 90.933 USD con las referencias del E-24. El escenario valoriza minutos-persona evitados y depende del volumen y del ahorro efectivo por pedido. Los beneficios con efectos superpuestos se evalúan individualmente. Este escenario de ahorro es independiente de la liquidación variable de la sección 1.2 (LafroX SpA, 2026a, secc.~2.7).

\subsection{Condiciones comerciales aplicables}
Las condiciones de ejecución se aplican a ambos alcances y vinculan la aceptación, los niveles de servicio y la gestión de cambios con las obligaciones del contrato.

\paragraph{Garantías y seguros.} La garantía de fiel cumplimiento asciende a USD 1.000.000 y se mantiene desde la firma hasta doce meses después del término de la operación. La garantía de correcto funcionamiento equivale al 5\,\% del valor total de implementación, dura 24 meses y es condición de la aceptación final del mes 21. LafroX mantiene y acredita las coberturas exigidas durante los períodos que fija el artículo 38 (Bases Administrativas, arts.~36--38, pp.~23--24).

\paragraph{Niveles de servicio.} El servicio gestionado se mide mediante los indicadores de disponibilidad, respuesta y resolución aplicables. LafroX entrega el informe mensual de servicio, con los incidentes, las causas y el plan de acción, dentro de los primeros cinco días hábiles. Las multas por incumplimiento se aplican conforme a las Bases y al E-25. La condición OTIF determina el devengo del componente INN-04 y se verifica por separado (Bases Administrativas, arts.~78--79, p.~40, y formulario E-25, p.~74; LafroX SpA, 2026b).

\paragraph{Control de cambios.} Toda modificación requiere solicitud formal, evaluación del impacto y aprobación previa del Comité Ejecutivo, además del acuerdo escrito de ambas partes. El límite acumulado es el 20\,\% del valor contractual original y se preservan el objeto, la naturaleza del contrato, las garantías mínimas y el cronograma obligatorio. Ejecutar un cambio sin aprobación previa no da derecho a pago adicional (Bases Administrativas, art.~72, p.~38).

\section*{Referencias}
\phantomsection
\addcontentsline{toc}{section}{Referencias}

\noindent Distribuidora Puelche S.A. (2026). \textit{Aclaraciones de la licitación} (Licitación TFEP-01/2026).

\noindent LafroX SpA. (2026a). \textit{Análisis económico-financiero} [Entregable 2 de la Oferta Económica].

\noindent LafroX SpA. (2026b). \textit{Innovaciones} [Subdocumento 13 y Formulario T-19 de la Oferta Técnica].

\noindent LafroX SpA. (2026c). \textit{Introducción al alcance de la solución} [Subdocumento 3 de la Oferta Técnica].

\noindent LafroX SpA. (2026d). \textit{Modelo financiero en planilla de cálculo} [Entregable 3 de la Oferta Económica, versión del 13 de noviembre de 2026].

\noindent LafroX SpA. (2026e). \textit{Plan de trabajo, EDT, cronograma e implantación} [Subdocumento 7 de la Oferta Técnica].

\section*{Declaración de uso de IA}
\phantomsection
\addcontentsline{toc}{section}{Declaración de uso de IA}

Se utilizó Codex para preparar y ampliar la redacción, ordenar la evidencia y cotejar cálculos y referencias. La Tabla~\ref{tab:e1-ia} declara el uso por sección. En esta versión de circulación interna, la revisión humana se registra cuando el equipo confirma el nombre de la persona revisora y qué comprobó.  (Bases Administrativas, Distribuidora Puelche S.A., 2026).

\begin{tablalafrox}{Uso de IA por sección}{tab:e1-ia}{F{0.10} F{0.12} F{0.32} F{0.10} F{0.12} F{0.24}}
{\cab{Sección} & \cab{Herramienta} & \cab{Finalidad del uso} & \cab{Nivel en texto} & \cab{Nivel en diagramas} & \cab{Revisión humana}}
1.1 & Codex & Redacción de condiciones de precio y vigencia, y cotejo de importes, conversiones e IVA. & Alto & Ninguno & Maximiliano Miño. Revisó la vigencia, las conversiones, los porcentajes y el IVA contra el E-21 y el E-24. \\
1.2 & Codex & Redacción de pagos, aceptación y liquidación variable, y cotejo del E-25 con el flujo de caja. & Alto & Ninguno & Tomas Perez. Verificó hitos, calendario, montos e IVA, y cotejó la liquidación variable con el E-25 y el modelo financiero. \\
1.3 & Codex & Redacción y síntesis del alcance, responsabilidades, beneficios y condiciones contractuales. & Alto & Ninguno & Maximiliano Miño. Revisó el alcance, las exclusiones, los beneficios y la coherencia de las condiciones comerciales. \\
\end{tablalafrox}
